\documentclass[11pt]{article}


\usepackage{fullpage}
\usepackage{graphicx}
\usepackage{amsmath}
\usepackage{amssymb}
\usepackage{amsthm}
\usepackage{fancyvrb}

\newcommand{\myname}{Mehshan Mustafa}

\newenvironment{theorem}[2][Theorem]{\begin{trivlist}
\item[\hskip \labelsep {\bfseries #1}\hskip \labelsep {\bfseries #2.}]}{\end{trivlist}}
\newenvironment{lemma}[2][Lemma]{\begin{trivlist}
\item[\hskip \labelsep {\bfseries #1}\hskip \labelsep {\bfseries #2.}]}{\end{trivlist}}
\newenvironment{exercise}[2][Exercise]{\begin{trivlist}
\item[\hskip \labelsep {\bfseries #1}\hskip \labelsep {\bfseries #2.}]}{\end{trivlist}}
\newenvironment{problem}[2][Problem]{\begin{trivlist}
\item[\hskip \labelsep {\bfseries #1}\hskip \labelsep {\bfseries #2.}]}{\end{trivlist}}
\newenvironment{question}[2][Question]{\begin{trivlist}
\item[\hskip \labelsep {\bfseries #1}\hskip \labelsep {\bfseries #2.}]}{\end{trivlist}}
\newenvironment{corollary}[2][Corollary]{\begin{trivlist}
\item[\hskip \labelsep {\bfseries #1}\hskip \labelsep {\bfseries #2.}]}{\end{trivlist}}
\newenvironment{solution}{\begin{proof}[Solution]}{\end{proof}}
\newenvironment{idea}[2][Proof Idea.]{\textit{#1} #2}



\parindent0in
\pagestyle{plain}
\thispagestyle{plain}

\usepackage{csquotes}
\usepackage[shortlabels]{enumitem}

\newcommand{\dated}{\today}

\begin{document}

\textbf{Introduction to the Theory of
Computation}\hfill\textbf{\myname}\\[0.01in]
\textbf{Chapter 1: Regular Languages}\hfill\textbf{\dated}\\
\smallskip\hrule\bigskip

\begin{problem}{1.63}
\end{problem}

\begin{problem}[Part]{1}
Let $A$ be an infinite regular language. Prove that $A$ can be split into two infinite disjoint regular subsets.
\end{problem}

\begin{proof}
If $A$ is a regular language, then according to Pumping lemma, some string $s$ in $A$ of length at least $p$ (the pumping length) can be divided into three pieces, $s = xyz$, where for any $i \geq 0$, the string $xy^iz \in A$. Split $A$ into two disjoint sets $A_1$ and $A_2$ as follows.

\[
A_1 = (A - \{s\}) \cup \{w \ | \ w = xy^iz, \text{ where } i \text{ is any odd integer greater than or equal to 0}\}. 
\]
\[
A_2 = (A - \{s\}) \cup \{w \ | \ w = xy^iz, \text{ where } i \text{ is any even integer greater than or equal to 0}\}. 
\]
\end{proof}

\begin{problem}[Part]{2}
Let $B$ and $D$ be two languages. Write $B \Subset D$ if $B \subseteq D$ and $D$ contains infinitely many strings that are not in $B$. Show that if $B$ and $D$ are two regular languages where $B \Subset D$, then we can find a regular language $C$ where $B \Subset C \Subset D$.
\end{problem}

\begin{proof}
\end{proof}

\end{document}